% --- Archon physics formalization source begin ---
% archon:physics
% archon:covers PhyXMiniProblems/problem_phyx_mini_0671.lean
% archon:source-report reports/phyx_mini/problem_phyx_mini_0671.source.json

\chapter{Physics problem phyx\_mini\_0671}
\label{ch:PhyXMiniProblems_problem_phyx_mini_0671}

\paragraph{Problem source.}
\#\# Physical scenario

A student drives a moped along a straight road as described by the velocity--time graph in figure

\#\# Question

What is the moped's final position at \textbackslash{}( t = 9.00 \textbackslash{}text\{ s\} \textbackslash{})?

\#\# Answer choices

- A: A: 16m

- B: B: 14m

- C: C: 28m

- D: D: 32m

\#\# Figure caption (auxiliary; use the image as primary evidence)

The image is a graph plotting velocity against time. The horizontal axis represents time \textbackslash{}( t \textbackslash{}) in seconds, labeled as \textbackslash{}( t \textbackslash{}, (\textbackslash{}text\{s\}) \textbackslash{}), ranging from 0 to 10 seconds. The vertical axis represents velocity \textbackslash{}( v\_x \textbackslash{}) in meters per second, labeled as \textbackslash{}( v\_x \textbackslash{}, (\textbackslash{}text\{m/s\}) \textbackslash{}), ranging from -8 to 8 m/s.

The plot is a piecewise linear graph:

1. From \textbackslash{}( t = 0 \textbackslash{}) to \textbackslash{}( t = 2 \textbackslash{}) seconds, the velocity increases linearly from 0 to 8 m/s.
2. From \textbackslash{}( t = 2 \textbackslash{}) to \textbackslash{}( t = 6 \textbackslash{}) seconds, the velocity remains constant at 8 m/s.
3. From \textbackslash{}( t = 6 \textbackslash{}) to \textbackslash{}( t = 10 \textbackslash{}) seconds, the velocity decreases linearly from 8 m/s to -8 m/s.

There is a grid in the background of the graph, aiding in better visualization of the data.

\#\# Dataset metadata

Kinematics; Spatial Relation Reasoning, Multi-Formula Reasoning

\paragraph{Recorded answer/context.}
C: C: 28m

\paragraph{Figure/image path.}
/root/proposal\_for\_physic/science-mango/phyx\_mini\_run/phyx\_data/test\_image/671.png

\paragraph{Formalization target.}
create a compiling Lean file with sorry bodies at `PhyXMiniProblems/problem\_phyx\_mini\_0671.lean`. The Lean declarations must preserve the physical quantities, dimensions or dimensional roles, figure labels, governing-law hypotheses, and final relation expressed by this problem.
Use Mathlib/Physlib names found through LeanExplore where available. If a physics API is missing, introduce faithful local abstractions rather than scalar placeholder aliases.

\begin{theorem}[Physics formalization target]
\label{thm:physics:phyx_mini_0671:target}
\lean{PhyXMiniProblems.ProblemPhyXMini0671.problem_phyx_mini_0671}
\uses{def:physics:phyx-mini-0671:phyxminiproblems-problemphyxmini0671-timefromseconds, def:physics:phyx-mini-0671:phyxminiproblems-problemphyxmini0671-positioninmeters, def:physics:phyx-mini-0671:phyxminiproblems-problemphyxmini0671-straightroadmopedsetup, def:physics:phyx-mini-0671:phyxminiproblems-problemphyxmini0671-matchesmopedscenario, def:physics:phyx-mini-0671:phyxminiproblems-problemphyxmini0671-matchesprimaryvelocitytimefigure, def:physics:phyx-mini-0671:phyxminiproblems-problemphyxmini0671-velocityreadoutmatchesfigure, def:physics:phyx-mini-0671:phyxminiproblems-problemphyxmini0671-startsatspatialorigin, def:physics:phyx-mini-0671:phyxminiproblems-problemphyxmini0671-satisfiespositionvelocitylaw, def:physics:phyx-mini-0671:phyxminiproblems-problemphyxmini0671-answerchoice, def:physics:phyx-mini-0671:phyxminiproblems-problemphyxmini0671-displayedfinalpositioninmeters}
The assigned autoformalize agent should translate this physics problem into Lean declarations in the covered file, with theorem and lemma proof bodies written as `by sorry`.
\end{theorem}
\begin{proof}
This is an autoformalization task, not a proof task. Produce statements that can later be proved by the physics prover without weakening the physical model.
\end{proof}

% --- Archon declaration topology begin ---
\section{Declaration topology}
\begin{definition}[Time Quantity]
  \label{def:physics:phyx-mini-0671:phyxminiproblems-problemphyxmini0671-timequantity}
  \lean{PhyXMiniProblems.ProblemPhyXMini0671.TimeQuantity}
  A unit-independent physical time coordinate
\end{definition}
\begin{proof}
  This is the declared model definition of Time Quantity.
\end{proof}
\begin{definition}[Position Quantity]
  \label{def:physics:phyx-mini-0671:phyxminiproblems-problemphyxmini0671-positionquantity}
  \lean{PhyXMiniProblems.ProblemPhyXMini0671.PositionQuantity}
  A signed physical position along the straight road
\end{definition}
\begin{proof}
  This is the declared model definition of Position Quantity.
\end{proof}
\begin{definition}[Signed Velocity Quantity]
  \label{def:physics:phyx-mini-0671:phyxminiproblems-problemphyxmini0671-signedvelocityquantity}
  \lean{PhyXMiniProblems.ProblemPhyXMini0671.SignedVelocityQuantity}
  A signed physical velocity component along the road's x-axis
\end{definition}
\begin{proof}
  This is the declared model definition of Signed Velocity Quantity.
\end{proof}
\begin{definition}[time Readout]
  \label{def:physics:phyx-mini-0671:phyxminiproblems-problemphyxmini0671-timereadout}
  \lean{PhyXMiniProblems.ProblemPhyXMini0671.timeReadout}
  \uses{def:physics:phyx-mini-0671:phyxminiproblems-problemphyxmini0671-timequantity}
  Numerical coordinate of a physical time in a selected time unit
\end{definition}
\begin{proof}
  This is the declared model definition of time Readout.
\end{proof}
\begin{definition}[position Readout]
  \label{def:physics:phyx-mini-0671:phyxminiproblems-problemphyxmini0671-positionreadout}
  \lean{PhyXMiniProblems.ProblemPhyXMini0671.positionReadout}
  \uses{def:physics:phyx-mini-0671:phyxminiproblems-problemphyxmini0671-positionquantity}
  Numerical coordinate of a signed physical position in a selected length unit
\end{definition}
\begin{proof}
  This is the declared model definition of position Readout.
\end{proof}
\begin{definition}[signed Velocity Readout]
  \label{def:physics:phyx-mini-0671:phyxminiproblems-problemphyxmini0671-signedvelocityreadout}
  \lean{PhyXMiniProblems.ProblemPhyXMini0671.signedVelocityReadout}
  \uses{def:physics:phyx-mini-0671:phyxminiproblems-problemphyxmini0671-signedvelocityquantity}
  Numerical signed-velocity component in selected length and time units
\end{definition}
\begin{proof}
  This is the declared model definition of signed Velocity Readout.
\end{proof}
\begin{definition}[time From Readout]
  \label{def:physics:phyx-mini-0671:phyxminiproblems-problemphyxmini0671-timefromreadout}
  \lean{PhyXMiniProblems.ProblemPhyXMini0671.timeFromReadout}
  \uses{def:physics:phyx-mini-0671:phyxminiproblems-problemphyxmini0671-timequantity}
  Construct a physical time coordinate from a readout in a selected unit
\end{definition}
\begin{proof}
  This is the declared model definition of time From Readout.
\end{proof}
\begin{definition}[time From Seconds]
  \label{def:physics:phyx-mini-0671:phyxminiproblems-problemphyxmini0671-timefromseconds}
  \lean{PhyXMiniProblems.ProblemPhyXMini0671.timeFromSeconds}
  \uses{def:physics:phyx-mini-0671:phyxminiproblems-problemphyxmini0671-timequantity, def:physics:phyx-mini-0671:phyxminiproblems-problemphyxmini0671-timefromreadout}
  SI-second time coordinate used by the graph
\end{definition}
\begin{proof}
  This is the declared model definition of time From Seconds.
\end{proof}
\begin{definition}[position In Meters]
  \label{def:physics:phyx-mini-0671:phyxminiproblems-problemphyxmini0671-positioninmeters}
  \lean{PhyXMiniProblems.ProblemPhyXMini0671.positionInMeters}
  \uses{def:physics:phyx-mini-0671:phyxminiproblems-problemphyxmini0671-positionquantity, def:physics:phyx-mini-0671:phyxminiproblems-problemphyxmini0671-positionreadout}
  SI-metre readout of a signed physical position
\end{definition}
\begin{proof}
  This is the declared model definition of position In Meters.
\end{proof}
\begin{definition}[signed Velocity In Meters Per Second]
  \label{def:physics:phyx-mini-0671:phyxminiproblems-problemphyxmini0671-signedvelocityinmeterspersecond}
  \lean{PhyXMiniProblems.ProblemPhyXMini0671.signedVelocityInMetersPerSecond}
  \uses{def:physics:phyx-mini-0671:phyxminiproblems-problemphyxmini0671-signedvelocityquantity, def:physics:phyx-mini-0671:phyxminiproblems-problemphyxmini0671-signedvelocityreadout}
  SI metres-per-second readout of a signed axial velocity
\end{definition}
\begin{proof}
  This is the declared model definition of signed Velocity In Meters Per Second.
\end{proof}
\begin{definition}[Driver Kind]
  \label{def:physics:phyx-mini-0671:phyxminiproblems-problemphyxmini0671-driverkind}
  \lean{PhyXMiniProblems.ProblemPhyXMini0671.DriverKind}
  Kind of driver stated in the scenario
\end{definition}
\begin{proof}
  This is the declared model definition of Driver Kind.
\end{proof}
\begin{definition}[Vehicle Kind]
  \label{def:physics:phyx-mini-0671:phyxminiproblems-problemphyxmini0671-vehiclekind}
  \lean{PhyXMiniProblems.ProblemPhyXMini0671.VehicleKind}
  Kind of road vehicle stated in the scenario
\end{definition}
\begin{proof}
  This is the declared model definition of Vehicle Kind.
\end{proof}
\begin{definition}[Road Geometry]
  \label{def:physics:phyx-mini-0671:phyxminiproblems-problemphyxmini0671-roadgeometry}
  \lean{PhyXMiniProblems.ProblemPhyXMini0671.RoadGeometry}
  Geometry of the road on which the vehicle moves
\end{definition}
\begin{proof}
  This is the declared model definition of Road Geometry.
\end{proof}
\begin{definition}[Motion Axis]
  \label{def:physics:phyx-mini-0671:phyxminiproblems-problemphyxmini0671-motionaxis}
  \lean{PhyXMiniProblems.ProblemPhyXMini0671.MotionAxis}
  The longitudinal spatial axis named by the velocity label vₓ
\end{definition}
\begin{proof}
  This is the declared model definition of Motion Axis.
\end{proof}
\begin{definition}[Figure Axis]
  \label{def:physics:phyx-mini-0671:phyxminiproblems-problemphyxmini0671-figureaxis}
  \lean{PhyXMiniProblems.ProblemPhyXMini0671.FigureAxis}
  The two labelled axes in the supplied plot
\end{definition}
\begin{proof}
  This is the declared model definition of Figure Axis.
\end{proof}
\begin{definition}[Curve Vertex]
  \label{def:physics:phyx-mini-0671:phyxminiproblems-problemphyxmini0671-curvevertex}
  \lean{PhyXMiniProblems.ProblemPhyXMini0671.CurveVertex}
  Named vertices of the piecewise-linear curve, in chronological order
\end{definition}
\begin{proof}
  This is the declared model definition of Curve Vertex.
\end{proof}
\begin{definition}[Velocity Time Point]
  \label{def:physics:phyx-mini-0671:phyxminiproblems-problemphyxmini0671-velocitytimepoint}
  \lean{PhyXMiniProblems.ProblemPhyXMini0671.VelocityTimePoint}
  A scalar coordinate read directly from the velocity--time graph
\end{definition}
\begin{proof}
  This is the declared model definition of Velocity Time Point.
\end{proof}
\begin{definition}[Velocity Time Figure]
  \label{def:physics:phyx-mini-0671:phyxminiproblems-problemphyxmini0671-velocitytimefigure}
  \lean{PhyXMiniProblems.ProblemPhyXMini0671.VelocityTimeFigure}
  \uses{def:physics:phyx-mini-0671:phyxminiproblems-problemphyxmini0671-figureaxis, def:physics:phyx-mini-0671:phyxminiproblems-problemphyxmini0671-curvevertex, def:physics:phyx-mini-0671:phyxminiproblems-problemphyxmini0671-velocitytimepoint}
  Labels, ranges, vertices, and presentation facts visible in image 671.png
\end{definition}
\begin{proof}
  This is the declared model definition of Velocity Time Figure.
\end{proof}
\begin{definition}[Straight Road Moped Setup]
  \label{def:physics:phyx-mini-0671:phyxminiproblems-problemphyxmini0671-straightroadmopedsetup}
  \lean{PhyXMiniProblems.ProblemPhyXMini0671.StraightRoadMopedSetup}
  \uses{def:physics:phyx-mini-0671:phyxminiproblems-problemphyxmini0671-timequantity, def:physics:phyx-mini-0671:phyxminiproblems-problemphyxmini0671-positionquantity, def:physics:phyx-mini-0671:phyxminiproblems-problemphyxmini0671-signedvelocityquantity, def:physics:phyx-mini-0671:phyxminiproblems-problemphyxmini0671-driverkind, def:physics:phyx-mini-0671:phyxminiproblems-problemphyxmini0671-vehiclekind, def:physics:phyx-mini-0671:phyxminiproblems-problemphyxmini0671-roadgeometry, def:physics:phyx-mini-0671:phyxminiproblems-problemphyxmini0671-motionaxis, def:physics:phyx-mini-0671:phyxminiproblems-problemphyxmini0671-velocitytimefigure}
  The physical motion. Position and velocity are independent dimensionful observables; neither is defined from a solved formula or an answer choice
\end{definition}
\begin{proof}
  This is the declared model definition of Straight Road Moped Setup.
\end{proof}
\begin{definition}[Matches Moped Scenario]
  \label{def:physics:phyx-mini-0671:phyxminiproblems-problemphyxmini0671-matchesmopedscenario}
  \lean{PhyXMiniProblems.ProblemPhyXMini0671.MatchesMopedScenario}
  \uses{def:physics:phyx-mini-0671:phyxminiproblems-problemphyxmini0671-straightroadmopedsetup}
  Categorical facts explicitly stated in the physical scenario
\end{definition}
\begin{proof}
  This is the declared model definition of Matches Moped Scenario.
\end{proof}
\begin{definition}[Matches Primary Velocity Time Figure]
  \label{def:physics:phyx-mini-0671:phyxminiproblems-problemphyxmini0671-matchesprimaryvelocitytimefigure}
  \lean{PhyXMiniProblems.ProblemPhyXMini0671.MatchesPrimaryVelocityTimeFigure}
  \uses{def:physics:phyx-mini-0671:phyxminiproblems-problemphyxmini0671-straightroadmopedsetup}
  Facts transcribed from the primary raster. In particular, the graph's actual breakpoints are 3 s and 5 s, and the drawn curve ends at 9 s, although the horizontal plot axis continues to 10 s
\end{definition}
\begin{proof}
  This is the declared model definition of Matches Primary Velocity Time Figure.
\end{proof}
\begin{definition}[Velocity Readout Matches Figure]
  \label{def:physics:phyx-mini-0671:phyxminiproblems-problemphyxmini0671-velocityreadoutmatchesfigure}
  \lean{PhyXMiniProblems.ProblemPhyXMini0671.VelocityReadoutMatchesFigure}
  \uses{def:physics:phyx-mini-0671:phyxminiproblems-problemphyxmini0671-timefromseconds, def:physics:phyx-mini-0671:phyxminiproblems-problemphyxmini0671-signedvelocityinmeterspersecond, def:physics:phyx-mini-0671:phyxminiproblems-problemphyxmini0671-straightroadmopedsetup}
  Interpretation of the three line segments as calibrated signed-velocity measurements. This contains the graph data but no displacement, position, or answer-choice value
\end{definition}
\begin{proof}
  This is the declared model definition of Velocity Readout Matches Figure.
\end{proof}
\begin{definition}[Starts At Spatial Origin]
  \label{def:physics:phyx-mini-0671:phyxminiproblems-problemphyxmini0671-startsatspatialorigin}
  \lean{PhyXMiniProblems.ProblemPhyXMini0671.StartsAtSpatialOrigin}
  \uses{def:physics:phyx-mini-0671:phyxminiproblems-problemphyxmini0671-timefromseconds, def:physics:phyx-mini-0671:phyxminiproblems-problemphyxmini0671-positioninmeters, def:physics:phyx-mini-0671:phyxminiproblems-problemphyxmini0671-straightroadmopedsetup}
  The implicit origin convention required for the question to have an absolute position answer rather than only a displacement answer
\end{definition}
\begin{proof}
  This is the declared model definition of Starts At Spatial Origin.
\end{proof}
\begin{definition}[Satisfies Position Velocity Law]
  \label{def:physics:phyx-mini-0671:phyxminiproblems-problemphyxmini0671-satisfiespositionvelocitylaw}
  \lean{PhyXMiniProblems.ProblemPhyXMini0671.SatisfiesPositionVelocityLaw}
  \uses{def:physics:phyx-mini-0671:phyxminiproblems-problemphyxmini0671-timefromseconds, def:physics:phyx-mini-0671:phyxminiproblems-problemphyxmini0671-positioninmeters, def:physics:phyx-mini-0671:phyxminiproblems-problemphyxmini0671-signedvelocityinmeterspersecond, def:physics:phyx-mini-0671:phyxminiproblems-problemphyxmini0671-straightroadmopedsetup}
  For every ordered pair of displayed time readouts, signed displacement equals the interval integral of signed velocity. This is a general physical law and does not contain the requested final position
\end{definition}
\begin{proof}
  This is the declared model definition of Satisfies Position Velocity Law.
\end{proof}
\begin{definition}[Answer Choice]
  \label{def:physics:phyx-mini-0671:phyxminiproblems-problemphyxmini0671-answerchoice}
  \lean{PhyXMiniProblems.ProblemPhyXMini0671.AnswerChoice}
  Labels of the four displayed multiple-choice answers
\end{definition}
\begin{proof}
  This is the declared model definition of Answer Choice.
\end{proof}
\begin{definition}[displayed Final Position In Meters]
  \label{def:physics:phyx-mini-0671:phyxminiproblems-problemphyxmini0671-displayedfinalpositioninmeters}
  \lean{PhyXMiniProblems.ProblemPhyXMini0671.displayedFinalPositionInMeters}
  \uses{def:physics:phyx-mini-0671:phyxminiproblems-problemphyxmini0671-answerchoice}
  Position in metres printed beside each displayed answer label
\end{definition}
\begin{proof}
  This is the declared model definition of displayed Final Position In Meters.
\end{proof}
\begin{definition}[recorded Dataset Answer]
  \label{def:physics:phyx-mini-0671:phyxminiproblems-problemphyxmini0671-recordeddatasetanswer}
  \lean{PhyXMiniProblems.ProblemPhyXMini0671.recordedDatasetAnswer}
  \uses{def:physics:phyx-mini-0671:phyxminiproblems-problemphyxmini0671-answerchoice}
  The answer label recorded by the source dataset; it is metadata, not a premise
\end{definition}
\begin{proof}
  This is the declared model definition of recorded Dataset Answer.
\end{proof}
\begin{lemma}[displacement at nine seconds]
  \label{lem:physics:phyx-mini-0671:phyxminiproblems-problemphyxmini0671-displacement-at-nine-seconds}
  \lean{PhyXMiniProblems.ProblemPhyXMini0671.displacement_at_nine_seconds}
  \uses{def:physics:phyx-mini-0671:phyxminiproblems-problemphyxmini0671-timefromseconds, def:physics:phyx-mini-0671:phyxminiproblems-problemphyxmini0671-positioninmeters, def:physics:phyx-mini-0671:phyxminiproblems-problemphyxmini0671-straightroadmopedsetup, def:physics:phyx-mini-0671:phyxminiproblems-problemphyxmini0671-velocityreadoutmatchesfigure, def:physics:phyx-mini-0671:phyxminiproblems-problemphyxmini0671-satisfiespositionvelocitylaw}
  The signed area under the graph through 9 s is 28 m: the initial triangle contributes 12, the plateau 16, and the positive and negative triangles under the final segment cancel
\end{lemma}
\begin{proof}
  The conclusion follows from the governing relations and figure readouts named in the statement of displacement at nine seconds.
\end{proof}
% --- Archon declaration topology end ---
% --- Archon physics formalization source end ---
